\begin{recipe}{Quiche Lorraine}
\kicker[Hoofdgerecht,Vlees,Frans,Hartig]{Hoofdgerecht · vlees · Frans · hartig}
\index[register]{Quiche Lorraine}
\index[register]{Ham}
\index[register]{Oude kaas}

\meta{1{,}5 uur \dvd Voor 6 porties \dvd Oven 180\,°C}

\lettrine{D}{e} klassieke quiche met ham en kaas. Deeg en werkwijze zijn hetzelfde
als bij de andere quiches in dit boek, alleen de vulling is eenvoudiger.

\blockrule{Ingrediënten}

\noindent\textit{Deeg}
\begin{ingredients}
  \ing{150 g}{bloem}
  \ing{75 g}{roomboter}
  \ing{1}{ei}
  \ing{1 tl}{zout}
\end{ingredients}

\medskip
\noindent\textit{Vulling}
\begin{ingredients}
  \ing{250 g}{hamblokjes}\index[register]{Ham}
  \ing{100 g}{oude kaas, geraspt}\index[register]{Oude kaas}
  \ing{4}{eieren}
  \ing{200 g}{crème fraîche}
  \ingb{versgemalen peper}
  \ingb{Bieslook voor garnering (optioneel)}
\end{ingredients}

\blockrule{Bereiding}
\tip{Haal de quiche uit de oven zodra het midden nog maar héél licht trilt, hij
gaart daarna vanzelf verder na. Te lang bakken maakt de vulling rubberachtig
in plaats van zijdezacht.}
\begin{steps}
  \item Kneed de deegingrediënten tot een samenhangend deeg. Maak een bol, wikkel
        in folie en laat 15 minuten rusten in de koelkast. Dat laat de boter weer
        koud worden en het gluten ontspannen, zodat het deeg zich makkelijker laat
        uitrollen en minder krimpt in de oven.
  \item Verwarm de oven voor op 180\,°C (boven- en onderwarmte). 
  \item Rol het deeg uit, bekleed de springvorm en
        prik met een vork wat gaatjes in de bodem: dat voorkomt dat het deeg
        tijdens het voorbakken opbolt. Bak de bodem 15 minuten voor.
  \item Bak de hamblokjes in een kleine koekenpan tot ze beginnen te kleuren. Haal van het
        vuur en laat iets afkoelen.
  \item Strooi wat kaas in de bodem en schep dan de ham in de voorgebakken bodem en strooi 
        dan de rest van de geraspte kaas erover.
  \item Klop de 4 eieren met de crème fraîche en wat peper tot een glad mengsel en
        giet dit erover.
  \item Bak de quiche 45 minuten op 180\,°C tot de vulling gestold is. 
  \item Laat de quiche 10 minuten rusten voor het aansnijden. Garneer eventueel met wat bieslook.
\end{steps}

\heroimagefade{images/quiche_loraine}
\end{recipe}
